\chapter{Conclusiones}
\label{cap:capitulo6}

Este TFM ha diseñado una propuesta de integración transversal de la sostenibilidad digital en el módulo MP0485 de DAM 1.º, con cinco prácticas (P0–P4) vinculadas a Resultados de Aprendizaje concretos, sustentadas en el ciclo de Kolb y evaluables mediante rúbricas analíticas. El análisis del BOE confirma que el currículo de Informática ofrece puntos de inserción naturales para trabajar huella de carbono y eficiencia energética sin modificar la programación oficial ni superar el 4,6 \% de la carga horaria del módulo. La revisión del estado del arte respalda que los enfoques más eficaces son aquellos que combinan medición empírica con datos propios del alumnado y aprendizaje basado en proyectos, que es el modelo que articula el arco P0–P4. Green IT no es un añadido externo al currículo de programación: es una dimensión que lo enriquece y le da relevancia profesional y ciudadana.


La principal limitación es que la propuesta no ha sido implementada en aula real y, por tanto, los indicadores de impacto propuestos tienen carácter orientativo. La generalización a otros ciclos o etapas requiere adaptar los RA y los entregables, y la medición del cambio de actitud a largo plazo queda fuera del alcance de este trabajo.


Como líneas futuras se proponen: un piloto en aula durante el curso 2026-27 con recogida de datos pre/post; la adaptación de la propuesta a otros módulos (ASIR, DAW) y a Bachillerato de Tecnología; y el desarrollo de un dashboard de aula que integre las métricas de CodeCarbon de todos los alumnos en tiempo real.
